\section{Evaluation Metrics}
\label{sec:metrics}

Let $y_i$ denote a reference GHI value, $\widehat y_i$ its forecast,
$e_i=y_i-\widehat y_i$, $\bar y$ the mean reference value, and $n$ the number
of evaluated targets. Metrics are calculated by location before the unweighted
macro-average is formed.

\subsection{Mean absolute error}

\begin{equation}
 \mathrm{MAE}=\frac{1}{n}\sum_{i=1}^{n}|e_i|.
 \label{eq:mae}
\end{equation}
MAE remains in W/m$^2$ and gives every absolute deviation linear weight. It is
easy to interpret but does not emphasize rare large errors.

\subsection{Root mean squared error}

\begin{equation}
 \mathrm{RMSE}=\sqrt{\frac{1}{n}\sum_{i=1}^{n}e_i^2}.
 \label{eq:rmse}
\end{equation}
RMSE is also expressed in W/m$^2$ and penalizes large deviations more strongly
than MAE. Its value can therefore be dominated by a small number of extreme
forecast errors.

\subsection{Normalized root mean squared error}

\begin{equation}
 \mathrm{nRMSE}(\%)=100\frac{\mathrm{RMSE}}{\bar y}.
 \label{eq:nrmse}
\end{equation}
nRMSE makes errors more comparable across locations with different mean GHI,
but becomes unstable when the denominator is close to zero. It is therefore
not computed on nighttime-only subsets.

\subsection{Coefficient of determination}

\begin{equation}
 R^2=1-\frac{\sum_{i=1}^{n}e_i^2}
 {\sum_{i=1}^{n}(y_i-\bar y)^2}.
 \label{eq:r2}
\end{equation}
$R^2$ measures tracked variation relative to a constant mean predictor. It can
be negative and should not replace scale-dependent error metrics.

The choice of a primary metric will be fixed for each resolution before the
final test comparison. MAE, RMSE, nRMSE, and $R^2$ provide complementary
evidence and will not be selectively reported according to which favors the
proposed model \cite{hyndman2006}.

